\section{Strategische Ma{\ss}nahmen im Berichtsjahr 2025}

Dieser Teil berichtet, was das Unternehmen im Gesch\"aftsjahr 2025 getan
hat, und bewertet es nicht. Wertungen stehen in Teil~\ref{sec:kurs} und
Teil~\ref{sec:urteil}.

\subsection{\"Ubernahme der Accession Risk Management Group}
\label{sec:accession}

Zum 1.\,August 2025 hat Brown \& Brown die Accession Risk Management Group
\"ubernommen. Die Gesamtgegenleistung betr\"agt
\MioUSDexakt{\AccessionKaufpreis}, davon \MioUSDexakt{\AccessionBarzahlung}
in bar und \MioUSDexakt{\AccessionAktien} in eigenen Aktien; der
Restbetrag entf\"allt auf eine sonstige Verbindlichkeit. Vom Kaufpreis
wurden \MioUSDexakt{\AccessionGoodwill} als Gesch\"afts- oder Firmenwert und
\MioUSDexakt{\AccessionImmat} als erworbene Kundenbeziehungen und sonstige
immaterielle Verm\"ogenswerte angesetzt; die gewichtete Nutzungsdauer der
Kundenbeziehungen betr\"agt \num{\NutzungsdauerKunden}
Jahre.\quelleGB{2025}{63}\quelleMerken{s8gb}

Es war nicht der einzige Erwerb des Jahres. Insgesamt hat das Unternehmen
2025 \num{\AkquisitionenAnzahl} Gesellschaften \"ubernommen und dabei
\num{\AkquisitionenKoepfe} Besch\"aftigte hinzugewonnen;\quelleGB{2025}{8}\quelleMerken{s8gb}
die Summe aller Gegenleistungen betr\"agt
\MioUSDexakt{\KaeufeKaufpreis}.\quelleGB{2025}{64}\quelleMerken{s8gb} In der
Gewinn- und Verlustrechnung 2025 sind daraus
\MioUSDexakt{\UmsatzAusKaeufen} Umsatz und
\MioUSDexakt{\ErgebnisAusKaeufen} Ergebnis vor Steuern
enthalten.\quelleErneut{s8gb}

Der Jahresumsatz der erworbenen Gesellschaften l\"asst sich aus zwei
ver\"offentlichten Angaben herleiten. Die Pro-forma-Rechnung, die alle
Zuk\"aufe des Jahres so behandelt, als w\"aren sie zu Beginn des Jahres 2024
erfolgt, weist f\"ur 2025 einen Umsatz von \MioUSDexakt{\ProFormaUmsatz}
aus, also \MioUSDexakt{\ProFormaMehrumsatz} mehr als
berichtet.\quelleErneut{s8gb} Rechnet man den bereits
konsolidierten Beitrag von \MioUSDexakt{\UmsatzAusKaeufen} hinzu, ergibt
sich ein Jahresumsatz der Zuk\"aufe von
\MioUSDexakt{\ZukaeufeUmsatzJahr} und damit ein Kaufpreis vom
\Faktor{\KaufpreisJeUmsatz}-fachen Jahresumsatz.

\subsection{Neuordnung der Segmente}

Im Zusammenhang mit der \"Ubernahme hat das Unternehmen seine Berichtsstruktur
von drei auf zwei Segmente umgestellt: Die vormaligen Segmente Programs und
Wholesale Brokerage wurden zur Specialty Distribution zusammengefasst,
neben der das Retail-Segment fortbesteht. Die Vorjahreszahlen wurden auf
die neue Struktur \"ubergeleitet, die Konzernabschl\"usse blieben davon
unber\"uhrt.\quelleGB{2025}{4}\quelleMerken{s9gb} Die Specialty Distribution gliedert sich in
Arrowhead Programs, Bridge Specialty Group und die neu aus dem
Accession-Teil One80 Intermediaries gebildete Arrowhead
Specialty.\quelleGB{2025}{42}\quelleMerken{s9gb}

\subsection{Finanzierung der \"Ubernahme}
\label{sec:finanzierung}

Die \"Ubernahme wurde im Juni 2025 in drei Schritten vorfinanziert. Erstens
begab das Unternehmen am 11.\,Juni sechs Anleihetranchen \"uber zusammen
\MioUSDexakt{\AnleihevolumenNeu} mit Kupons zwischen \Pct{\KuponMin} und
\Pct{\KuponMax} und Laufzeiten von 2026 bis 2055.\quelleGB{2025}{45}\quelleMerken{s9gb}
Zweitens platzierte es am 10.\,Juni \num{\KapitalerhoehungStueck} neue
Aktien zu \USDje{\KapitalerhoehungKurs} und erl\"oste daraus netto
\MioUSDexakt{\KapitalerhoehungNetto}. Drittens gingen Aktien im Wert von
\MioUSDexakt{\VerkaeuferaktienWert} unmittelbar an die verkaufenden
Gesellschafter, bemessen zum Schlusskurs von
\USDje{\VerkaeuferaktienKurs} am 6.\,Juni 2025.\quelleGB{2025}{46}\quelleMerken{s9gb}

Die gesamten Finanzverbindlichkeiten stiegen dadurch von
\MioUSDexakt{\GesamtverschuldungVJ} auf
\MioUSDexakt{\Gesamtverschuldung}, die Anleihen allein von
\MioUSDexakt{\AnleihenVJ} auf \MioUSDexakt{\Anleihen}.\quelleGB{2025}{67}\quelleMerken{s9gb}

\subsection{Kapitalallokation neben der \"Ubernahme}

Am 22.\,Oktober 2025 hat das Board den Rahmen f\"ur Aktienr\"uckk\"aufe um
\MioUSDexakt{\RueckkaufAufstockung} aufgestockt; zum Jahresende standen davon noch
\MioUSDexakt{\RueckkaufRahmen} offen. Im Jahr 2025 selbst wurden
\num{\RueckkaufStueck} Aktien zu einem Durchschnittskurs von
\USDje{\RueckkaufKurs} und damit f\"ur \MioUSDexakt{\AktienrueckkaufCF}
zur\"uckgekauft.\quelleGB{2025}{28}\quelleMerken{s9gb} Im ersten Halbjahr 2026 hat das
Unternehmen den R\"uckkauf deutlich ausgeweitet: Der Bestand eigener Aktien
zu Anschaffungskosten stieg von \MioUSDexakt{\HJEigeneAktienVJ} auf
\MioUSDexakt{\HJEigeneAktien}, also um
\MioUSDexakt{\HJEigeneAktienZuwachs}.\quelleQM{Form 10-Q Q2 2026}{7}\quelleMerken{s9qmformqq}

\subsection{Wirkung auf das Ergebnis}

Der Umsatz stieg 2025 um \Pct{\UmsatzDelta}, das bereinigte EBITDAC um
\Pct{\EBITDACberDelta}, das Konzernergebnis um
\Pct{\ErgebnisDelta} und das verw\"asserte Ergebnis je Aktie um
\Pct{\EPSdilDelta}.\quelleGB{2025}{49}\quelleMerken{s9gb} Die Kette wird von Stufe zu Stufe
flacher. Zwischen EBITDAC und Konzernergebnis liegen der auf
\MioUSDexakt{\Zinsaufwand} gestiegene Zinsaufwand und die auf
\MioUSDexakt{\AbschrImmat} gestiegene Abschreibung auf erworbene
immaterielle Verm\"ogenswerte; zwischen Konzernergebnis und Ergebnis je
Aktie liegt der Zuwachs der Aktienzahl um \Pct{\AktienZuwachs}.

Segmentseitig tr\"agt die Specialty Distribution die St\"uckmarge: Ihre
bereinigte EBITDAC-Marge stieg von \Pct{\fpeval{round(\SpezMargeVJ,1)}} auf
\Pct{\fpeval{round(\SpezMarge,1)}}, w\"ahrend die des Retail-Segments mit
\Pct{\fpeval{round(\RetailMarge,1)}} unver\"andert blieb.\quelleGB{2025}{39}\quelleMerken{s9gb}

\subsection{Ausblick und dessen Nachpr\"ufbarkeit}
\label{sec:ausblick}

Als US-amerikanischer Emittent ver\"offentlicht Brown \& Brown im
Jahresbericht keine Prognose in Form einer Umsatz- oder Margenbandbreite.
Ein Abgleich von Ausblick und tats\"achlichem Ergebnis, wie ihn ein
deutscher Gesch\"aftsbericht erm\"oglicht, ist deshalb nicht durchf\"uhrbar
und wird hier nicht ersatzweise konstruiert.

Nachpr\"ufbar ist stattdessen der Verlauf des Fr\"uhindikators aus
Abschnitt~\ref{sec:organisch}. Er hat sich nach dem Bilanzstichtag weiter
verschlechtert: Im ersten Halbjahr 2026 betrug das organische Wachstum
\Pct{\HJOrganisch} nach \Pct{\HJOrganischVJ} im Vorjahreszeitraum, im
zweiten Quartal allein \Pct{\QzweiOrganisch} nach
\Pct{\QzweiOrganischVJ}.\quelleQM{Form 10-Q Q2 2026}{29}\quelleMerken{s10qmformqq}
Der berichtete Umsatz stieg im selben Halbjahr um \Pct{\HJUmsatzDelta} auf
\MioUSDexakt{\HJUmsatz}, getragen von der ganzj\"ahrigen Konsolidierung der
Zuk\"aufe.\quelleErneut{s10qmformqq}
